\documentclass{article}
\usepackage[UTF8]{ctex}
\usepackage{amsmath,amssymb,amsfonts}
\usepackage{enumitem}
\usepackage{geometry}
\geometry{a4paper, margin=1in}
% === 错题本排版优化（不参与题目标识） ===
\setlength{\parindent}{0pt}
\setlength{\parskip}{0.35em plus 0.10em minus 0.08em}
\setlength{\emergencystretch}{3em}
\setlist{itemsep=0.25em, topsep=0.40em, parsep=0pt}
\allowdisplaybreaks[2]
% === 错题本排版优化结束 ===
\begin{document}
\section*{河北工业大学2024年高等代数真题}
\section*{58．河北工业大学2024年高等代数试卷}

\subsection{题目1}
一．（15 分）设 \(f_{0}(x), f_{1}(x), \cdots, f_{n-1}(x)\) 为 \(n\) 个实系数多项式，且

\[
\left(x^{n}-2\right) \mid \sum_{i=0}^{n-1} f_{i}\left(x_{i}^{n}\right) x^{i} .
\]

证明：\(x-2 \mid f_{i}(x), i=0,1, \cdots, n-1\) ．

\subsection{题目2}
二．（15 分）计算行列式

\[
\left|\begin{array}{cccccc}
x & 0 & 0 & \cdots & 0 & 1 \\
-1 & x & 0 & \cdots & 0 & 2 \\
0 & -1 & x & \cdots & 0 & 3 \\
\vdots & \vdots & \vdots & & \vdots & \vdots \\
0 & 0 & 0 & \cdots & x & 2023 \\
0 & 0 & 0 & \cdots & -1 & x+2024
\end{array}\right| .
\]

\subsection{题目3}
三．（15 分）设 \(A\) 为 5 阶矩阵，\(k\) 为大于 5 的正整数，证明：\(r\left(A^{5}\right)=r\left(A^{k}\right)\) 。

\subsection{题目4}
四．（15 分）设 \(A\) 为4阶实对称矩阵，且 \(A\) 的特征值为 \(\lambda_{1}=\lambda_{2}=1, \lambda_{3}=\lambda_{4}=2\) ，且 \((1,1,0,0)^{\prime},(1,1,0\)。为 \(A\) 的属于特征值1的特征向量，求矩阵 \(A\) 。

\subsection{题目5}
五．（15 分）设线性方程组的一个解为 \(\eta\) ，其导出组的一个基础解系为 \(\xi_{1}, \xi_{2}, \cdots, \xi_{s}\) ，记

\[
\gamma_{0}=\eta, \gamma_{1}=\eta+\xi_{1}, \gamma_{2}=\eta+\xi_{2}, \cdots, \gamma_{s}=\eta+\xi_{s} .
\]

证明：方程组的任意一个解 \(\gamma\) 可表示为 \(\sum_{i=0}^{s} \mu_{i} \gamma_{i}\) ，其中 \(\sum_{i=0}^{s} \mu_{i}=1\) 。

\subsection{题目6}
六．（15 分）设 \(W\) 是数域 \(P\) 上所有形如

\[
\left(\begin{array}{ccccc}
a_{1} & a_{2} & a_{3} & \cdots & a_{n} \\
a_{n} & a_{1} & a_{2} & \cdots & a_{n-1} \\
a_{n-1} & a_{n} & a_{1} & \cdots & a_{n-2} \\
\vdots & \vdots & \vdots & \vdots & \vdots \\
a_{2} & a_{3} & a_{4} & \cdots & a_{1}
\end{array}\right)
\]

的矩阵关于矩阵的加法和数量乘法构成的线性空间，\(U\) 是数域 \(P\) 上 \(n\) 阶对角阵全体关于矩阵的加法并数量乘法构成的线性空间，证明：\(W\) 和 \(U\) 同构。

\subsection{题目7}
七．（20 分）设 \(V_{1}\) 是数域 \(P\) 上全体 \(n\) 阶对称矩阵的集合，\(V_{2}\) 是数域 \(P\) 上全体 \(n\) 阶反称矩阵的集合，\(P^{n \times r}\) 是 \(n\) 阶矩阵生成的线性空间。

\begin{itemize}
  \item[（1）] 证明：\(P^{\mathrm{n} \times \mathrm{n}}=V_{1} \oplus V_{2}\) ．
  \item[（2）] 证明：存在唯一的线性变换 \(\sigma\) ，使得 \(\sigma\) 的值域为 \(V_{1}\) ，核为 \(V_{2}\) ，且 \(\sigma^{2}=\sigma\) 。
\end{itemize}

\subsection{题目8}
八．（20分）设 \(A, B\) 均为 \(n\) 阶实对称矩阵，且 \(B\) 。为正定矩阵。

\begin{itemize}
  \item[（1）] 若 \(A\) 为半正定矩阵，证明：\(A+B\) 为正定矩阵。
  \item[（2）] 证明：存在一实可逆矩阵 \(S\) ，使得 \(S^{\prime} A B\) 与 \(S^{\prime} B S\) 均为对角形．
\end{itemize}

\subsection{题目9}
九．（20 分）设 \(A\) 为任意 \(m \times n\) 矩阵，\(X\) 为未知 \(n \times m\) 矩阵．证明：矩阵方程 \(A X A=A\) 一定有解．
\end{document}
